When a reader processes a passage, many signals about \emph{how} the text was produced --- who wrote it, whether a human or a language model wrote it, in what register, and even what keyboard layout shaped its typos --- coexist in the same string of tokens. We ask whether these \emph{production modes} share a unified representational substrate inside a mid-size language model, and whether they are geometrically disentangled. We probe four very different modes in the last-token hidden states of \qwen: 7-way literary authorship, AI-vs-human authorship, formal-vs-casual register, and QWERTY-vs-Dvorak typo patterns. A single linear probe decodes every mode well above chance, with peak accuracies of 79.2\%--99.8\% across the four modes. The probe directions are pairwise near-orthogonal (all $|\cos| \leq 0.18$), and cross-mode transfer of probe directions is near chance for orthogonal modes (\typo $\leftrightarrow$ \register at 0.48) while exposing one expected content overlap (\register $\to$ \aihuman at 0.74, reflecting LLM formality bias). Layer-depth profiles split the four modes into three distinct tiers: content/register modes saturate at shallow layers, authorship emerges gradually across the stack, and surface-typing artifacts become decodable only at the final layers. Production modes are real, multiple, and largely independent latent variables in LLM hidden states --- a unification with consequences for authorship privacy, controllable generation, and AI-text detection.
